% =====================================================================
% Apresentacao do autor — compartilhada pelos volumes de estudo
% Manter este texto alinhado com a apresentacao da apostila principal.
% =====================================================================
\section*{Apresentação}
\addcontentsline{toc}{section}{Apresentação}
\thispagestyle{fancy}

Este material foi desenvolvido como ferramenta prática de apoio ao ensino e à
formação técnica de estudantes e profissionais da área elétrica, especialmente
daqueles envolvidos com \textbf{Eletricidade Básica}, \textbf{Análise de
Circuitos Elétricos} e \textbf{Eletromagnetismo}.

A proposta desta coletânea é oferecer uma base sólida e progressiva de
conhecimentos por meio da resolução de problemas contextualizados, que auxiliam
na fixação dos conceitos teóricos e no desenvolvimento do raciocínio aplicado à
eletricidade e ao magnetismo.

\begin{conceito}[Como usar este material]
\textbf{1. Níveis de dificuldade.} Cada questão traz uma marcação de nível:
\nivelmarca{1}~\emph{fixação} (aplicação direta de uma definição),
\nivelmarca{2}~\emph{aplicação} (exige interpretação e uma ou duas etapas de
cálculo), \nivelmarca{3}~\emph{aprofundamento} (várias etapas ou combinação de
conceitos) e \nivelmarca{4}~\emph{desafio} (síntese, demonstração, otimização ou
modelagem não rotineira). O Nível 4 é usado apenas quando a complexidade realmente
justifica essa classificação. Dentro de cada seção as questões estão
\textbf{ordenadas em ordem crescente de dificuldade}.

\smallskip
\textbf{2. Gabarito comentado.} Ao final de cada seção há a chave de respostas
seguida das \emph{resoluções comentadas}. Resolva primeiro; consulte depois.

\smallskip
\textbf{3. Lembretes teóricos.} As caixas azuis no início de cada seção reúnem as
fórmulas e os princípios necessários para resolver as questões daquele bloco.
\end{conceito}

\vspace{1em}
\begin{flushright}
\textit{Desejo a todos uma excelente jornada de estudos.}\\[2mm]
\textit{Prof. Francisco}
\end{flushright}
